\section{LLM Collaboration Workflow}

The MICO collaboration loop has five stages:

\begin{enumerate}
  \item inventory modules, ports, protocols, domains, and adapters;
  \item ask the model for a structured graph or JSON AST;
  \item run compiler checks and return structured diagnostics;
  \item ask for a minimal JSON patch when repair is enabled;
  \item run RTL validation and emit traceability from checked artifacts.
\end{enumerate}

This separates proposal from validation. The LLM never becomes the source of correctness; its output is constrained to data structures that the compiler can reject, repair, and compare across model profiles.

The provider policy separates cost from correctness. Smoke tests and prompt development use low-cost profiles first; higher-cost profiles are reserved for candidates that pass structured-output, compiler, and RTL gates.

The provider harness records validate-only and authenticated runs as redacted JSON artifacts: profile, model identifier, prompt hash, request settings, JSON validity, token/cost metadata when available, repair turns, compiler attachment, and EDA attachment. It records whether an API key was present but never stores or prints the key. The batch runner applies this to Direct Verilog, SystemVerilog-interface, MICO source, MICO JSON AST, and JSON AST repair baselines over the benchmark manifest.
